% modules/open-targets/A1-data-informed-glm.tex
% -----------------------------------------------------------------------------
% Deep-dive companion to the A1 entry of the Part I agenda (modules/08-open-targets)
% and to Question q:glm-data-informed of the GLM spine (modules/05-tier3-glm):
% data-informed functional-inequality constants for Gaussian-prior GLM posteriors,
% via the Brascamp--Lieb / inverse-Hessian route plus a certified posterior-tail
% correction.
% Compiles inside main.tex; standalone it shows cross-module \ref as "??".
% -----------------------------------------------------------------------------
\documentclass[../../main.tex]{subfiles}
\begin{document}
\section{A1 --- Data-informed constants for GLM posteriors}
\label{sec:a1}

This is the deep-dive companion to the A1 paragraph of the agenda (Section~\ref{sec:agenda}) and
to Question~\ref{q:glm-data-informed} of the GLM spine (Section~\ref{sec:tier3}). It is the
flagship finite-sample target: replace the data-free certificate
$\CP\le\CLS\le\lmax(\Sigmaz)$ of Theorem~\ref{thm:glm-fi}, which reads only the prior curvature,
by a bound that tracks the observed information, of the advertised shape
\begin{equation}\label{eq:a1-slogan}
  \CP(\pi)\ \lesssim\ \lmax\!\Big[\big(\Sigmaz^{-1}+X^\top\bar W X\big)^{-1}\Big]
  \ +\ \text{tail correction},
\end{equation}
interpolating the prior-dominated and likelihood-dominated regimes. The purpose of this section is
to make A1 precise, to split it into a Poincar\'e half that is essentially within reach of existing
machinery and a log-Sobolev/transport half that is genuinely harder, to record exactly which
inequality already supplies \eqref{eq:a1-slogan} and what remains to be done, and to state the
theorem one should actually aim to prove. In slogan form: for \(\CP\), the functional-inequality
skeleton and a certified bounded-curvature implementation now exist; the remaining work is to
prove a checkable finite-sample improvement over the prior and to generalize beyond that class.

\paragraph{Reviewed status.}
Proposition~\ref{prop:a1-euclidean-harmonic}, Proposition~\ref{prop:a1-mode-leverage}, and
Corollary~\ref{cor:a1-leverage-asymptotic} have passed independent review and are proved under
their displayed bounded-curvature hypotheses. Proposition~\ref{prop:a1-bulk-tail} and the
inverse-mode-Hessian no-go calculation below remain useful analytic drafts pending separate
review; neither is promoted by the reviewed bounded-curvature dossier.

\subsection{Setup: the observed-information geometry of a GLM posterior}
\label{subsec:a1-setup}

Fix the Gaussian-prior convex GLM posterior of Family~\ref{fam:glm-posterior},
\[
  \pi(\dd\theta)\ \propto\ e^{-U(\theta)}\dd\theta,
  \qquad
  U(\theta)\ =\ \tfrac12(\theta-m)^\top\Sigmaz^{-1}(\theta-m)\ +\ \sum_{i=1}^n\ell_i(x_i^\top\theta),
\]
with each $\ell_i(\cdot)=\ell(y_i,\cdot)$ convex and $C^2$ in the linear predictor and prior mean
$m$ (the agenda took $m=0$; nothing below depends on it). Writing $X=(x_i^\top)_{i=1}^n$ and the
diagonal weight matrix $W(\theta)=\diag\big(\ell_i''(x_i^\top\theta)\big)_{i=1}^n\succeq0$, the
Hessian is the observed-information geometry of the model,
\begin{equation}\label{eq:a1-hessian}
  H(\theta)\ :=\ \Hess U(\theta)\ =\ \Sigmaz^{-1}\ +\ X^\top W(\theta)X .
\end{equation}
For logistic regression $\ell_i''(t)=\sigma(t)(1-\sigma(t))=\tfrac14\operatorname{sech}^2(t/2)$,
the variance weight that vanishes as the fitted probability approaches $0$ or $1$; for Poisson
$\ell_i''(t)=e^{t}$; for the Gaussian/linear model $W\equiv\Id$. The prior-only Bakry--\'Emery
reading of Theorem~\ref{thm:glm-fi} keeps only $H(\theta)\succeq\Sigmaz^{-1}$ and discards the
PSD likelihood term, giving the safe but data-blind $\CP\le\lmax(\Sigmaz)$. The statistical
target is that, in the posterior bulk, the relevant curvature is not $\Sigmaz^{-1}$ but the full
\begin{equation}\label{eq:a1-target-curv}
  A_{\bar W}\ :=\ \Sigmaz^{-1}+X^\top\bar W X,
\end{equation}
for a deterministic, data-informed $\bar W\succeq0$ --- the nonasymptotic functional-inequality
analogue of the asymptotic GLM principle ``posterior covariance $\approx$ inverse observed
information'' that A2 (Section~\ref{sec:a2}, Conjecture~\ref{conj:a2}) makes exact as $n\to\infty$.

\subsection{The key existing result: Brascamp--Lieb gives \texorpdfstring{$\CP$}{CP} for free}
\label{subsec:a1-bl}

The Poincar\'e half of A1 is not literature-free: for log-concave measures there is a known
bridge from \emph{local} Hessian information to a \emph{global} Poincar\'e bound, and it is
exactly the tool \eqref{eq:a1-slogan} wants. Brascamp--Lieb (Theorem~\ref{thm:brascamp-lieb})
gives, for the strictly convex $U$ of \eqref{eq:a1-hessian},
\begin{equation}\label{eq:a1-bl}
  \Var_\pi(f)\ \le\ \E_\pi\,\inner{H(\theta)^{-1}\nabla f}{\nabla f}
  \ \le\ \E_\pi\!\big[\lmax\!\big(H(\theta)^{-1}\big)\,\norm{\nabla f}^2\big] ,
\end{equation}
the second step bounding the local quadratic form pointwise by its largest eigenvalue. This is
\emph{not} yet a Poincar\'e inequality: the curvature factor $\lmax(H^{-1})$ and the gradient
$\norm{\nabla f}^2$ remain \emph{coupled} inside a single expectation, and one may not simply
factor $\E_\pi[\lmax(H^{-1})\norm{\nabla f}^2]$ into
$\big(\E_\pi\lmax(H^{-1})\big)\,\E_\pi\norm{\nabla f}^2$ --- that would require a negative
correlation between the two factors that need not hold. The passage to a genuine, unweighted
Poincar\'e inequality is not a direct decoupling of the weighted Dirichlet form but a two-step
self-improvement for log-concave measures: Brascamp--Lieb first yields a spread/variance bound
for $1$-Lipschitz $f$ (the average inverse curvature controls the fluctuations of Lipschitz
observables), and Milman's equivalence of concentration and spectral gap on log-concave measures
then upgrades that weak control to a full spectral gap. This is exactly the
Milman spread-to-gap route of Cattiaux and Guillin
\cite{cattiaux2018poincare}, which yields a genuine dimension-free bound of the form
\begin{equation}\label{eq:a1-cfg}
  \CP(\pi)\ \lesssim\ \E_\pi\,\lmax\!\big(H(\theta)^{-1}\big)
  \ =\ \E_\pi\,\lmax\!\Big[\big(\Sigmaz^{-1}+X^\top W(\theta)X\big)^{-1}\Big].
\end{equation}
This is the abstract inequality one would want to specialize to GLM posteriors. The crucial point
is that the operative quantity is $\E_\pi\,\lmax(H^{-1})$, the average \emph{inverse} curvature,
\emph{not} $\lmax[(\E_\pi H)^{-1}]$: matrix inversion is operator-convex, so by Jensen
$\E_\pi H^{-1}\succeq(\E_\pi H)^{-1}$, and large average curvature does not by itself imply small
average inverse curvature. The tail behaviour of $W(\theta)$ is therefore not a technical
annoyance but structurally part of the problem --- a single posterior direction along which
$W(\theta)\to0$ can dominate the expectation in \eqref{eq:a1-cfg}.

\subsection{The bulk--tail decomposition}
\label{subsec:a1-bulktail}

The data-informed bound follows by splitting the posterior into a high-curvature bulk and a
low-curvature tail. The universal Milman comparison culminating in
Proposition~\ref{prop:a1-bulk-tail} is an analytic proof draft and remains outside the reviewed set.
The open GLM content begins when one asks for a useful, computable choice of $\bar W$.
For a deterministic diagonal $\bar W\succeq0$ define the \emph{good
bulk-curvature event}
\begin{equation}\label{eq:a1-good}
  G_{\bar W}\ :=\ \big\{\theta:\ X^\top W(\theta)X\ \succeq\ X^\top\bar W X\big\}.
\end{equation}
On $G_{\bar W}$ one has $H(\theta)\succeq A_{\bar W}$, hence $H(\theta)^{-1}\preceq A_{\bar
W}^{-1}$; off $G_{\bar W}$ the Gaussian prior still gives $H(\theta)\succeq\Sigmaz^{-1}$, hence
$H(\theta)^{-1}\preceq\Sigmaz$. Splitting the expectation in \eqref{eq:a1-cfg} across $G_{\bar W}$
and its complement,
\begin{equation}\label{eq:a1-split}
  \E_\pi\,\lmax\!\big(H(\theta)^{-1}\big)
  \ \le\ \lmax\!\big(A_{\bar W}^{-1}\big)
  \ +\ \E_\pi\!\big[\lmax\!\big(H(\theta)^{-1}\big);\,G_{\bar W}^c\big]
  \ \le\ \lmax\!\big(A_{\bar W}^{-1}\big)+\lmax(\Sigmaz)\,\pi(G_{\bar W}^c),
\end{equation}
where the last inequality is the crude probabilistic form of the tail correction. Combining with
\eqref{eq:a1-cfg} gives precisely the A1 shape \eqref{eq:a1-slogan}:
\begin{equation}\label{eq:a1-main}
  \CP(\pi)\ \lesssim\
  \underbrace{\lmax\!\Big[\big(\Sigmaz^{-1}+X^\top\bar W X\big)^{-1}\Big]}_{\text{bulk Fisher curvature}}
  \ +\
  \underbrace{\lmax(\Sigmaz)\,\pi(G_{\bar W}^c)}_{\text{penalty for low-curvature mass}} .
\end{equation}
The sharper and more honest version keeps the integrated inverse-curvature tail rather than
factoring out $\lmax(\Sigmaz)$,
\begin{equation}\label{eq:a1-main-sharp}
  \CP(\pi)\ \lesssim\ \lmax\!\Big[\big(\Sigmaz^{-1}+X^\top\bar W X\big)^{-1}\Big]
  \ +\ \E_\pi\!\big[\lmax\!\big(H(\theta)^{-1}\big);\,G_{\bar W}^c\big],
\end{equation}
with \eqref{eq:a1-main} the clean corollary. This already interpolates the two regimes: in the
likelihood-dominated bulk the first term is at the inverse-observed-information scale
$\lmax[(\Sigmaz^{-1}+X^\top\bar W X)^{-1}]\ll\lmax(\Sigmaz)$, and the second term is small exactly
when the data-informed $\bar W$ certifies the bulk curvature on a high-probability set.

\begin{proposition}[Universal inverse-Hessian bulk--tail certificate]
\label{prop:a1-bulk-tail}
There is a universal $C_M<\infty$ such that every posterior in
Subsection~\ref{subsec:a1-setup} and every deterministic diagonal $\bar W\succeq0$ satisfy
\eqref{eq:a1-main-sharp} and \eqref{eq:a1-main}, with the implicit constant equal to $C_M$.
\end{proposition}

\begin{proof}
For every $1$-Lipschitz $f$, Brascamp--Lieb gives
$\Var_\pi(f)\le\E_\pi\lmax(H^{-1})$. Milman's log-concave spread-to-gap comparison, in the
form recorded by \cite{cattiaux2018poincare}, therefore yields
$\CP(\pi)\le C_M\E_\pi\lmax(H^{-1})$. The two Loewner bounds used in
\eqref{eq:a1-split} finish the proof.
\end{proof}

\begin{proposition}[Factor-one Euclidean harmonic-mean bound]
\label{prop:a1-euclidean-harmonic}
Let $U\in C^\infty(\R^d)$ and suppose that, for finite $0<m\le M$,
\[
  mI\preceq\Hess U(\theta)\preceq MI\qquad(\theta\in\R^d).
\]
Writing $\rho(\theta)=\lmin(\Hess U(\theta))$, one has
\begin{equation}\label{eq:a1-harmonic-one}
  \CP(\pi)\ \le\ \E_\pi\frac1{\rho(\theta)}
  \ =\ \E_\pi\lmax\!\big((\Hess U(\theta))^{-1}\big).
\end{equation}
In particular, \eqref{eq:a1-main-sharp} and \eqref{eq:a1-main} hold with factor one for every
finite Gaussian-prior binary-logistic posterior.
\end{proposition}

\begin{proof}
Let $A=-\Delta+\nabla U\cdot\nabla$ be the nonnegative self-adjoint generator on $L^2(\pi)$ and
$\lambda_1=1/\CP(\pi)$. For $\varepsilon>0$, choose a nonzero $f$ in the spectral subspace of
$A$ over $[\lambda_1,\lambda_1+\varepsilon]$, and put $g=\norm{\nabla f}$. Spectral calculus,
the integrated weighted Bochner identity, and the weak Kato inequality give
\[
 (\lambda_1+\varepsilon)\E_\pi g^2
 \ \ge\ \E_\pi\norm{\nabla g}^2+\E_\pi[\rho g^2].
\]
Applying the Poincar\'e inequality with gap $\lambda_1$ to $g$ yields
\[
 \E_\pi[\rho g^2]\le
 \lambda_1(\E_\pi g)^2+\varepsilon\E_\pi g^2.
\]
Cauchy--Schwarz gives
$(\E_\pi g)^2\le\E_\pi[\rho g^2]\E_\pi[1/\rho]$, while
$\E_\pi g^2\le m^{-1}\E_\pi[\rho g^2]$. Dividing and sending
$\varepsilon\downarrow0$ proves \eqref{eq:a1-harmonic-one}.

For completeness, the noncompact domain step follows from the two-sided Hessian hypothesis, not
from the compact source theorem. Conjugating $A$ by $e^{-U/2}$ gives the Schr\"odinger operator
$-\Delta+\frac14\norm{\nabla U}^2-\frac12\Delta U$. Its potential is bounded below and tends
quadratically to infinity because $\norm{\nabla U(\theta)}\ge m\norm\theta-O(1)$ and
$\Delta U\le dM$. Thus $C_c^\infty$ is an operator core. The compact-support Bochner identity
extends in the graph norm; positivity of $\Hess U$ controls the Hessian term, and the regularization
$(\norm{\nabla f}^2+\delta^2)^{1/2}$ closes the Kato step. For logistic regression,
$\Sigmaz^{-1}\preceq H(\theta)\preceq\Sigmaz^{-1}+\frac14X^\top X$, so the hypotheses hold.
\end{proof}

\begin{remark}[What factor one does not yet cover]
The compact theorem of \cite{Veysseire2011SpectralGap} motivates
\eqref{eq:a1-harmonic-one}, but the proof above supplies a separate noncompact Euclidean domain
argument. It does not establish the result for arbitrary strongly log-concave potentials with
unbounded Hessian, including the Poisson GLM, nor for nonsmooth or merely weakly convex losses.
Closing the Bochner identity on the spectral domain in those classes, or proving the necessary
Mosco and uniform-integrability limits under approximation, remains open. The universal Milman
factor $C_M$ in Proposition~\ref{prop:a1-bulk-tail} remains the available general draft route.
\end{remark}

\subsection{What the saturating-likelihood tail obstruction actually rules out}
\label{subsec:a1-tail}

For saturating likelihoods such as logistic regression, a positive mode or bulk likelihood
curvature cannot be promoted to a global curvature floor without controlling where that curvature
decays. This is the narrow flat-direction obstruction of
Proposition~\ref{prop:flat-prior-nonintegrable} read spectrally. For nearly separable logistic data
there are directions $v$
along which many linear predictors $x_i^\top(\theta+tv)\to\pm\infty$ as $t\to\infty$, so
\[
  w_i(\theta+tv)\ =\ \sigma\big(x_i^\top(\theta+tv)\big)\big(1-\sigma(x_i^\top(\theta+tv))\big)\ \to\ 0,
\]
and $X^\top W(\theta+tv)X$ loses rank in the limit. Hence no bound $X^\top W(\theta)X\succeq
B\succ0$ can hold for all $\theta$: the prior keeps the posterior proper and supplies tail drift,
but it cannot convert vanishing logistic curvature into a global likelihood floor.

There is also an analytic no-go family for replacing the whole posterior geometry by the Hessian
at the mode. In one dimension let
\[
  \pi_a(\dd\theta)\ \propto\
  e^{-\theta^2/(2\sigma^2)}\,\operatorname{sigmoid}(a\theta)\,\dd\theta,
  \qquad a\to\infty.
\]
The density relative to $N(0,\sigma^2)$ is $2\operatorname{sigmoid}(a\theta)$, so dominated
convergence gives a half-Gaussian limit and
$\Var_{\pi_a}(\theta)\to\sigma^2(1-2/\pi)>0$. If
$s_a=a\hat\theta_a$ at the mode, the score equation is
$s_a(1+e^{s_a})=a^2\sigma^2$, and hence
\[
 \big(\sigma^{-2}+a^2\operatorname{sigmoid}(s_a)
       (1-\operatorname{sigmoid}(s_a))\big)^{-1}
 =\frac{\sigma^2}{1+s_a\operatorname{sigmoid}(s_a)}\longrightarrow0.
\]
The linear test therefore makes the ratio of $\CP(\pi_a)$ to this inverse mode Hessian diverge.
This is an analytic proof draft of a precise no-go statement: no universal constant-times
inverse-mode-Hessian bound holds for this logistic family.

For a bulk certificate based on nonzero data curvature $\bar W$, the honest correction
consequently has the shape
\[
  \text{bulk Fisher curvature}\ +\ \text{penalty for posterior mass where that curvature fails},
\]
which is exactly the structure of \eqref{eq:a1-main}. The separable / weakly-identified logistic
regime is the stress test: a good theorem should expose how the constant degrades with the
separation geometry and the prior scale $\lmax(\Sigmaz)$, not pretend the degradation away.
This conclusion is deliberately narrower than ``every data-informed bound needs an explicit tail
term'': Gaussian linear models have $W\equiv\Id$, the prior-only bound is tail-free, and
covariance- or full-distribution-based bounds are not ruled out.

\subsection{High-dimensional anisotropy: the matrix bound does the right thing}
\label{subsec:a1-highdim}

The bulk term $\lmax[(\Sigmaz^{-1}+X^\top\bar W X)^{-1}]$ has the correct behaviour when $D>n$ or
$X^\top\bar W X$ is rank-deficient, which a scalar floor $\rho\Id$ cannot express. In directions
$v$ invisible to the data, $Xv=0$, so
$v^\top(\Sigmaz^{-1}+X^\top\bar W X)v=v^\top\Sigmaz^{-1}v$ and the bound reverts to the prior
scale; in directions the data inform, the curvature improves to the observed-information scale.
The single matrix expression thus interpolates
\[
  \begin{aligned}
    \text{data-dominated directions}&\ \leadsto\ \text{inverse-Fisher scale},\\
    \text{prior-only directions}&\ \leadsto\ \text{prior-covariance scale}.
  \end{aligned}
\]
which is exactly what a Bayesian GLM constant should do, and is the anisotropic refinement of the
isotropic floor $\lmax(\Sigmaz)$ of Theorem~\ref{thm:glm-fi}.

\subsection{Choosing \texorpdfstring{$\bar W$}{Wbar}: where the real work is}
\label{subsec:a1-barw}

The functional-inequality skeleton is the easy part; the technical content of A1 is the
construction of a good deterministic, data-informed $\bar W$ together with a certified bound on
$\pi(G_{\bar W}^c)$. There are three natural levels, of increasing sharpness.

\emph{(i) Mode-centred rowwise weights.} For logistic regression $w_i(t)$ is even and decreasing
in $\abs t$, so on a mode-centred ellipsoid $E_r=\{\theta:\norm{\theta-\hat\theta}_{A_{\bar
W}}\le r\}$ with $\abs{x_i^\top(\theta-\hat\theta)}\le\delta_i$ one has the explicit floor
\begin{equation}\label{eq:a1-rowwise}
  \bar w_i\ =\ \inf_{\abs{t-x_i^\top\hat\theta}\le\delta_i}w_i(t)
  \ =\ \frac{1}{4\cosh^2\!\big((\abs{x_i^\top\hat\theta}+\delta_i)/2\big)},
\end{equation}
and $G_{\bar W}^c\subseteq E_r^c$, reducing the tail to the posterior probability of leaving a
mode-centred ellipsoid. Rowwise floors are simple but conservative when $n$ is large.

\emph{(ii) Spectral / quantile weights.} The right object is the matrix inequality
\eqref{eq:a1-good}, not the rowwise $w_i(\theta)\ge\bar w_i$ for every $i$. One should choose
$\bar W$ to certify $X^\top W(\theta)X\succeq X^\top\bar W X$ with high posterior probability ---
e.g.\ as a posterior quantile of $\lmin[A_{\bar W}^{-1/2}X^\top W(\theta)X\,A_{\bar W}^{-1/2}]$ ---
which is far less conservative than requiring every weight to clear its floor.

\emph{(iii) Certifying the tail probability $\pi(G_{\bar W}^c)\le\eps$.} This is the heart of the
problem, and admits a hierarchy of certificates: prior-only concentration (safe but crude),
Brascamp--Lieb covariance control \eqref{eq:a1-bl} (data-informed second moments), and
finite-sample Laplace / posterior-concentration bounds (sharpest, model-specific)
\cite{KasprzakGiordanoBroderick2025}. An attractive self-improving scheme bootstraps these: the
prior-only $\CP\le\lmax(\Sigmaz)$ gives a crude bulk probability, which sharpens the constant via
\eqref{eq:a1-main}, which in turn tightens the bulk probability --- the iterative flavour of the
variable-curvature literature \cite{cattiaux2018poincare,CattiauxGuillin2022Perturbed}. The
constrained-Gibbs Poincar\'e inequalities of \cite{ConstrainedGibbs2026}, which already prove
$O(1/n)$ local constants on a high-probability good set for logistic and Poisson models, are the
natural starting toolkit, alongside the local log-Polyak--{\L}ojasiewicz bounds of
\cite{GongHeShen2025LogPL} and the Lyapunov-potential criteria of \cite{ChenSridharan2024Lyapunov}.

\subsection{An explicit mode-leverage certificate}
\label{subsec:a1-mode-leverage}

For an important GLM class the rowwise construction and its posterior tail can in fact be made
fully deterministic. Assume
\begin{equation}\label{eq:a1-loghess-lipschitz}
  \ell_i''(s)>0,
  \qquad
  \abs{\log\ell_i''(s)-\log\ell_i''(t)}\le L_i\abs{s-t}.
\end{equation}
Logistic and Poisson regression satisfy \eqref{eq:a1-loghess-lipschitz} with $L_i=1$, while the
Gaussian linear model has $L_i=0$. Let $\hat\theta$ be the posterior mode, put
\begin{equation}\label{eq:a1-leverage}
 \widehat H=H(\hat\theta),\qquad
 \alpha_i=L_i\sqrt{x_i^\top\widehat H^{-1}x_i},\qquad
 \eta=\max_i\alpha_i,
\end{equation}
and write $z=\widehat H^{1/2}(\theta-\hat\theta)$ and $r=\norm z$.

Define, with their continuous values $g_\pm(r;0)=r^2/2$,
\begin{equation}\label{eq:a1-gpm}
 g_-(r;\eta)=\frac{e^{-\eta r}+\eta r-1}{\eta^2},
 \qquad
 g_+(r;\eta)=\frac{e^{\eta r}-\eta r-1}{\eta^2}.
\end{equation}
\begin{proposition}[Deterministic mode-leverage certificate]
\label{prop:a1-mode-leverage}
Under \eqref{eq:a1-loghess-lipschitz},
\begin{equation}\label{eq:a1-hessian-sandwich}
 e^{-\eta r}I\preceq
 \widehat H^{-1/2}H(\theta)\widehat H^{-1/2}
 \preceq e^{\eta r}I
\end{equation}
and
\begin{equation}\label{eq:a1-potential-sandwich}
 g_-(r;\eta)\le U(\theta)-U(\hat\theta)\le g_+(r;\eta).
\end{equation}
For $R>0$, choose
\begin{equation}\label{eq:a1-leverage-wbar}
 \bar w_i(R)=e^{-\alpha_iR}\ell_i''(x_i^\top\hat\theta),
 \qquad E_R=\{\theta:\norm{\widehat H^{1/2}(\theta-\hat\theta)}\le R\}.
\end{equation}
Then $E_R\subseteq G_{\bar W(R)}$ and
\begin{equation}\label{eq:a1-leverage-tail}
 \pi(G_{\bar W(R)}^c)\le\pi(E_R^c)\le
 \varepsilon_d(R,\eta):=\min\!\left\{1,
 \frac{\displaystyle\int_R^\infty r^{d-1}e^{-g_-(r;\eta)}\,dr}
      {\displaystyle\int_0^\infty r^{d-1}e^{-g_+(r;\eta)}\,dr}\right\}.
\end{equation}
Thus, conditional on the separately unreviewed universal comparison in
Proposition~\ref{prop:a1-bulk-tail}, equations~\eqref{eq:a1-leverage-wbar}--
\eqref{eq:a1-leverage-tail} give a certificate computable from $(X,y,\Sigmaz)$ and
one-dimensional quadrature. The direct bounded-Hessian conclusion below does not use that
universal comparison.

If, in addition, Proposition~\ref{prop:a1-euclidean-harmonic} applies and $0\le\eta<1$, then
\begin{equation}\label{eq:a1-leverage-direct}
 \CP(\pi)\le K_d(\eta)\lmax(\widehat H^{-1}),
 \qquad
 K_d(\eta):=
 \frac{\displaystyle\int_0^\infty r^{d-1}e^{\eta r-g_-(r;\eta)}\,dr}
      {\displaystyle\int_0^\infty r^{d-1}e^{-g_+(r;\eta)}\,dr}.
\end{equation}
For fixed $d$, $K_d(\eta)\to1$ as $\eta\downarrow0$.
\end{proposition}

\begin{proof}
From \eqref{eq:a1-loghess-lipschitz},
$e^{-L_i\abs{x_i^\top(\theta-\hat\theta)}}\ell_i''(x_i^\top\hat\theta)
\le\ell_i''(x_i^\top\theta)\le
e^{L_i\abs{x_i^\top(\theta-\hat\theta)}}\ell_i''(x_i^\top\hat\theta)$.
Cauchy--Schwarz in the $\widehat H$ geometry bounds the exponent by $\alpha_i r\le\eta r$.
The fixed prior Hessian lies between its own $e^{-\eta r}$ and $e^{\eta r}$ multiples, so summing
the likelihood and prior terms proves \eqref{eq:a1-hessian-sandwich}. Integrating that comparison
twice along the segment from the mode proves \eqref{eq:a1-potential-sandwich}.

On $E_R$, the individual curvature floors in \eqref{eq:a1-leverage-wbar} hold, proving
$E_R\subseteq G_{\bar W(R)}$. In $z$ coordinates, the numerator of $\pi(E_R^c)$ is at most the
radial integral involving $g_-$, while its normalizing denominator is at least the radial integral
involving $g_+$; the sphere area and whitening Jacobian cancel. This proves
\eqref{eq:a1-leverage-tail}. Finally, \eqref{eq:a1-hessian-sandwich} gives
$\lmax(H(\theta)^{-1})\le e^{\eta r}\lmax(\widehat H^{-1})$.
Equation~\eqref{eq:a1-harmonic-one} and the same radial comparison prove
\eqref{eq:a1-leverage-direct}. Its numerator is finite for $\eta<1$, and dominated convergence
gives $K_d(\eta)\to1$.
\end{proof}

\begin{corollary}[Vanishing-leverage Poincar\'e limit]
\label{cor:a1-leverage-asymptotic}
Fix $d$. For a sequence of bounded-Hessian posteriors satisfying
\eqref{eq:a1-loghess-lipschitz}, if $\eta_n\to0$, then
\begin{equation}\label{eq:a1-leverage-asymptotic}
 \CP(\pi_n)=(1+o(1))\lmax(\widehat H_n^{-1}).
\end{equation}
If the inputs are random, $\eta_n\to0$ and $\widehat H_n/n\to I(\theta_0)\succ0$ in probability,
then $n\CP(\pi_n)\to\lmax(I(\theta_0)^{-1})$ in probability.
\end{corollary}

\begin{proof}
The potential sandwich and dominated convergence show that the $\widehat H_n$-whitened posterior
converges to $N(0,I_d)$ in total variation and second moments. Explicitly, fix
$0<\eta_0<1$. For all sufficiently large $n$ the unnormalized density is bounded by
$e^{-g_-(r;\eta_0)}$; this has an exponential tail, so its product with $1+r^2$ is integrable.
Meanwhile the two radial normalizing integrals squeeze the normalization to the Gaussian one.
This supplies the required uniform integrability. Hence the whitened covariance is $I_d+o(1)$,
and the linear test supplies
$\CP(\pi_n)\ge(1-o(1))\lmax(\widehat H_n^{-1})$. The matching upper bound is
\eqref{eq:a1-leverage-direct} with $K_d(\eta_n)\to1$.
\end{proof}

The condition $\eta<1$ is a genuine certificate gate, not a claim that the posterior lacks a
Poincar\'e inequality when it fails. In the one-observation family of
Subsection~\ref{subsec:a1-tail}, $\eta_a^2=a^2/H(\hat\theta_a)\to\infty$, so the direct certificate
correctly refuses the false inverse-mode-Hessian conclusion. Formula~\eqref{eq:a1-leverage-tail}
and the prior bound remain available. The exact corollary currently covers finite logistic
posteriors; Poisson satisfies the radial hypothesis but has unbounded Hessian, so factor one and
the exact coefficient require the unresolved domain extension noted above.

\subsection{Poincar\'e has a bounded-curvature theorem; log-Sobolev and transport are harder}
\label{subsec:a1-pls}

The three constants are \emph{not} on the same footing, and A1 should not pretend they are.

\emph{Poincar\'e.} The route of Subsections~\ref{subsec:a1-bl}--\ref{subsec:a1-bulktail} is
essentially complete in outline: Brascamp--Lieb \eqref{eq:a1-bl} plus the log-concave
self-improvement \cite{cattiaux2018poincare} gives \eqref{eq:a1-cfg}, and the bulk--tail split
gives \eqref{eq:a1-main}. Proposition~\ref{prop:a1-mode-leverage} supplies a computable $\bar W$
and a certified radial tail for the log-Hessian-Lipschitz class. The remaining work is quantitative
comparison with the prior baseline and less conservative spectral constructions for high-leverage
or unbounded-Hessian instances.

\emph{Log-Sobolev.} Here the naive ``replace $\Sigmaz^{-1}$ by $A_{\bar W}$'' is most misleading.
Bakry--\'Emery gives the prior-only $\CLS\le\lmax(\Sigmaz)$ (Theorem~\ref{thm:glm-fi}), but a Brascamp--Lieb-style
integrated-inverse-Hessian estimate does \emph{not} imply an LSI: log-Sobolev is sensitive to
global landscape and tails, not only to local curvature (the low-temperature companion result of
\cite{ChewiStromme2024Ballistic} identifies the ballistic LSI limit with the global
Polyak--{\L}ojasiewicz constant, which can be strictly worse than the local curvature; cf.\ A2,
Subsection on the LSI half). The relevant tools are Lyapunov criteria, defective LSI, and
super-Poincar\'e inequalities \cite{CattiauxGuillinWangWu2009,BakryCattiauxGuillin2008,CattiauxGuillin2010}.
A plausible data-informed target is a local LSI on the bulk $G_{\bar W}$ glued to Gaussian-prior
drift outside, of the shape
\begin{equation}\label{eq:a1-lsi-target}
  \CLS(\pi)\ \lesssim\ \lmax\!\big(A_{\bar W}^{-1}\big)\cdot\Phi\big(\pi(G_{\bar W}^c)\big)
  \ +\ (\text{prior-tail drift term}),
\end{equation}
with $\Phi$ possibly logarithmic --- but one should not claim the clean bound \eqref{eq:a1-main}
for $\CLS$ without additional global hypotheses.

For binary logistic likelihoods there is a stronger obstruction. Because the likelihood is
nonnegative, convex, and grows at most linearly, its Gaussian-prior posterior satisfies the exact
global identity
\begin{equation}\label{eq:a1-logistic-rigidity}
  \CLS(\pi)=\CTCI(\pi)=\lmax(\Sigmaz)
\end{equation}
for every finite data set; see Proposition~\ref{prop:a2-logistic-global}. Thus no global
bulk-improved LSI or $T_2$ theorem is possible for the flagship logistic model. The meaningful
second stage is explicitly local, restricted, warm-start, or uses a modified metric.

\emph{Transport.} The Talagrand $T_2$ constant inherits the LSI story: whenever a data-informed
LSI holds, Otto--Villani (Theorem~\ref{thm:otto-villani}) gives $\CTCI\le\CLS$ at the same scale,
while $T_2(C)$ itself implies $\Cov_\pi\preceq C I$ through its linear sub-Gaussian consequence.
A data-informed $\CTCI$ therefore follows a data-informed $\CLS$, not the
Poincar\'e bound alone. The variational-inference companion A4 (Section~\ref{sec:a4},
Question~\ref{q:a4-certificate}) pursues exactly this $\Wtwo$ analogue at posterior scale.

\subsection{The certified bounded-curvature result and the residual target}
\label{subsec:a1-theorem}

The abstract Poincar\'e calculation is Proposition~\ref{prop:a1-bulk-tail}, while
Proposition~\ref{prop:a1-mode-leverage} supplies a concrete selection theorem for a nontrivial
model class. The latter and its factor-one input are independently audited; the former universal
Milman comparison remains an analytic draft pending separate review. The honest positive target
is now to quantify the certified construction, prove that it beats the prior on a checkable
finite-sample regime, and improve it when maximal row leverage is too conservative.

\begin{conjecture}[Strict improvement on data-informative GLM subclasses]
\label{conj:a1}
Fix $d$ and a desired improvement margin $0<\delta<1$. There should be explicit, checkable
data-informativeness conditions on $(X,y,\Sigmaz)$ --- excluding, in particular, $X=0$ and
likelihood-flat prior-dominated instances --- under which the bounded-Hessian construction of
Proposition~\ref{prop:a1-mode-leverage} selects a nonempty range of $R$ for which
\[
  \begin{aligned}
    B_{\rm cert}(R)
    &:=\Big\{\lmax\!\Big[\big(\Sigmaz^{-1}+X^\top\bar W(R)X\big)^{-1}\Big]
      +\E_\pi\!\big[\lmax\!\big(H(\theta)^{-1}\big);\,G_{\bar W(R)}^c\big]\Big\}\\
    &\le\Big\{\lmax\!\big(A_{\bar W(R)}^{-1}\big)
      +\lmax(\Sigmaz)\,\varepsilon_d(R,\eta)\Big\}
  \end{aligned}
\]
satisfies
$\CP(\pi)\le B_{\rm cert}(R)\le(1-\delta)\lmax(\Sigmaz)$. The conditions must be expressed in
computable mode curvature, leverage, and radial-tail quantities, not in posterior-oracle
quantiles. For $\eta<1$, the already proved direct certificate gives the immediate sufficient
check
$K_d(\eta)\lmax(\widehat H^{-1})<\lmax(\Sigmaz)$; the residual content is to obtain useful split
or spectral certificates when that conservative check fails. Along regular fixed-dimensional
logistic sequences, Corollary~\ref{cor:a1-leverage-asymptotic} already gives the exact A2
coefficient.
For unbounded-Hessian GLMs, the separate analytic Milman route retains an explicit universal
factor and the exact coefficient remains open. The global $\CLS$/$\CTCI$ analogue is excluded for finite
Gaussian-prior logistic regression by \eqref{eq:a1-logistic-rigidity}; local or restricted
analogues remain open.
\end{conjecture}

\noindent Propositions~\ref{prop:a1-euclidean-harmonic} and~\ref{prop:a1-mode-leverage} now supply
the independently audited bounded-Hessian functional-inequality and selection skeletons. The
remaining program is to quantify when the explicit expression improves on the prior and extend or
replace maximal row leverage in the regimes where it fails. The universal Milman route remains
separate and pending review: its fixed comparison loss cannot by itself recover an exact
asymptotic coefficient.

\subsection{Sub-targets and what remains open}
\label{subsec:a1-open}

The flagship question (Question~\ref{q:glm-data-informed}) now splits into one prerequisite
constant/source-applicability task and five open selection, tail, sharpness, entropy, and
application tasks.

\begin{question}[Beyond bounded Euclidean curvature]
\label{q:a1-poincare}
Determine whether the spectral-domain argument of
Proposition~\ref{prop:a1-euclidean-harmonic} extends beyond globally bounded Hessian to
unbounded-Hessian posteriors such as Poisson, and to nonsmooth or only weakly convex GLMs. For the
general fallback, independently audit and make explicit the universal constant in
Proposition~\ref{prop:a1-bulk-tail}. The compact source theorem
\cite{Veysseire2011SpectralGap} alone does not discharge these noncompact domain questions.
\end{question}

\begin{question}[Computable, data-informed $\bar W$]
\label{q:a1-barw}
Use the explicit mode-leverage weights \eqref{eq:a1-leverage-wbar} as the deterministic baseline,
and compare them with matrix-spectral alternatives that do not let a single high-leverage row
control every direction. Any replacement must still be computable from $(X,y,\Sigmaz)$ alone and
must certify $X^\top W(\theta)X\succeq X^\top\bar W X$ without posterior-oracle quantiles.
\end{question}

\begin{question}[Sharpening the posterior-tail probability]
\label{q:a1-tail}
Quantify and optimize the finite-sample loss in the certified radial bound
\eqref{eq:a1-leverage-tail}, and find sharper matrix-tail bounds when maximal leverage is
conservative.
For GLMs outside \eqref{eq:a1-loghess-lipschitz}, prove an alternative computable tail certificate
via finite-sample Laplace bounds \cite{KasprzakGiordanoBroderick2025} or a self-improving bootstrap.
\end{question}

\begin{question}[Sharp constants and the $1/n$ scale]
\label{q:a1-sharp}
Verify the hypotheses of Corollary~\ref{cor:a1-leverage-asymptotic} for regular random-design
logistic sequences. For unbounded-Hessian models or sequences without vanishing maximal leverage,
find the correct replacement spectral-stability condition. Every claim must match the A2 limit
\eqref{eq:a2-target} and the lower bound
$\CP\ge\lmax(\Cov_\pi)$ (Remark~\ref{rem:rayleigh-lower-only}).
\end{question}

\begin{question}[Log-Sobolev and transport analogues]
\label{q:a1-lsi}
In view of \eqref{eq:a1-logistic-rigidity}, formulate and prove a \emph{localized or restricted}
data-informed LSI/transport inequality on the bulk $G_{\bar W}$, with an explicit admissible class
of initial laws or test functions and a tail/warm-start penalty. Determine which stronger-tail
GLMs admit a genuine global version. Lyapunov, super-Poincar\'e, and decomposition methods
\cite{CattiauxGuillinWangWu2009,BakryCattiauxGuillin2008,CattiauxGuillin2010,MenzSchlichting2014}
are the relevant tools; a global posterior-scale claim for ordinary logistic regression is not.
\end{question}

\begin{question}[Certified sampling and concentration rates]
\label{q:a1-sampling}
Translate the data-informed constants into Langevin / MALA complexity and posterior-concentration
guarantees, where $\CP,\CLS$ control convergence
\cite{ChewiEtAl2024LMC,GadatPanloupPellegrini2020}. This converts statistical curvature into
certified computational Bayesian rates, and is the payoff that motivates the whole program.
\end{question}

\begin{remark}[Relation to the rest of the agenda]
\label{rem:a1-relations}
A1 is the finite-sample flagship; A2 (Section~\ref{sec:a2}) is its $n\to\infty$ shadow, and a
sharp A1 bound tight in the bulk should reproduce the Fisher-information limit
\eqref{eq:a2-target}. The transport face of A1 is pursued by A4 (Section~\ref{sec:a4},
Question~\ref{q:a4-certificate}) as restricted/localized $\Wtwo$ constants at the same
$\lmax[(\Sigmaz^{-1}+X^\top\bar W X)^{-1}]$ bulk scale. The tail correction of \eqref{eq:a1-main}
is the bulk-side shadow of the heavy-tailed obstruction of A3 (Section~\ref{sec:a3}), where the
same vanishing-curvature directions survive into the posterior \emph{tail} and force weighted
constants; and the flat, weakly identified directions that inflate that tail term recur as the
symmetry- and funnel-induced slow modes that A5 (Section~\ref{sec:a5}) separates and quotients
away. The general convex-geometry boundary ---
$\CP\le K\lmax(\Cov_\pi)$ for \emph{every} log-concave measure with $K$ universal --- is the
Kannan--Lov\'asz--Simonovits conjecture of Part~III; A1 is the structured, finite-sum counterpart
that the KLS machinery does not exploit, which is precisely what makes it more tractable.
\end{remark}
\end{document}
